\documentclass[10pt,a4paper]{amsart}
\usepackage[margin=19mm]{geometry}
\usepackage{amsmath,amssymb,mathtools}
\usepackage[scr=rsfs,cal=euler]{mathalfa}
\usepackage{xfrac}
\usepackage[unicode,hidelinks,bookmarks=false]{hyperref}
\newcommand{\ie}{i.e.\ }
\newcommand{\cf}{cf.\ }
\newcommand{\resp}{respectively\ }
\newcommand{\Subsection}{\subsection}
\providecommand{\nobreakdash}{\nobreak\hbox{-}}
\setlength{\emergencystretch}{2em}
\multlinegap0pt
\usepackage{amssymb}
\newtheorem{thm}{Theorem}
\newcommand{\Pten}{\mathbf P}
\newcommand{\kten}{\mathbf k}
\title{The Penrose conjecture for initial data sets satisfying a $2$-convexity condition}
\author{Conghan Dong}
\date{Source: arXiv:2607.20741v2 (CC BY 4.0)}
\begin{document}
\maketitle

\noindent\emph{Source and licence.} The Penrose conjecture for initial data sets satisfying a $2$-convexity condition. Version arXiv:2607.20741v2 (CC BY 4.0), distributed by arXiv under the Creative Commons Attribution 4.0 International licence, http://creativecommons.org/licenses/by/4.0/, as recorded in the arXiv licence metadata for this exact version and shown on the abstract page https://arxiv.org/abs/2607.20741v2. Full licence notice: \texttt{licenses/arxiv-2607.20741-CC-BY-4.0.txt}.\par\smallskip

\noindent\textit{Editorial scope.} Counted unit: the article's thm thm-weak-mono (source0.tex lines 459--461). The article's complete original proof of it is delivered with it (source0.tex lines 463--494, one \texttt{proof} environment of 32 lines). No mathematics is written, repaired or re-derived: the statement and the proof are the article's own lines, in the article's own notation, and no retained line is substituted or re-flowed.  Retained uncounted as the source-local context the proof needs: lines 68--70; lines 407--414; lines 441--448.  Nothing was omitted from any retained span, so the package carries no omission record.  The retained lines cite 1 works, whose entries are transcribed verbatim from the article's own bibliography source in the e-print and delivered with short editorial labels recorded in the item's reference mapping.  No retained line contains a figure environment, an image-inclusion command or drawing code, so no image is delivered and the package declares no asset dependency.  Editorial formatting only, in addition: an AMS \texttt{amsart} preamble that restates the article's own declarations and macros used by the retained lines, namely the environments \texttt{thm} on separate counters, together with the standard packages those lines need.  The retained environments print in this document as Theorem 1 in order of appearance, whereas the article prints the same items with the numbering of its own full document.  The complete native source of the article as distributed in the arXiv e-print for this exact version is bundled unchanged as \texttt{sources/205830/source0.tex}.
\begin{align}\label{eq:DEC}
  \mu\geq |J|_g.
\end{align}

\begin{align}\label{nabla-P-id}
	\begin{split}
		(&\nabla _{\nu } \Pten) (\nu ,\nu ) \\
		&= (\mathrm{div}_{g} \Pten) (\nu ) - \mathrm{div}_{N_t}(\Pten ^{TN}(\nu ,\cdot ) ) -H \mathbf{P}(\nu ,\nu ) + \langle \mathrm{II}, \mathbf{P}^{TN} \rangle \\
&= (\mathrm{div}_{g} \Pten) (\nu ) - \mathrm{div}_{N_t}(\Pten ^{TN}(\nu ,\cdot ) )
						   + \frac{1}{2} H\mathrm{tr}_g \Pten - \frac{3}{2}H \Pten(\nu ,\nu )   + \langle \mathring{\Pten} ^{TN}, \mathring{\mathrm{II}} \rangle,
	\end{split}
\end{align} 

\begin{align}\label{B'-ineq}
	\begin{split}
		16 \pi & e ^{-\frac{t}{2}}B'(t) \\
	&\geq  \int_{N_t} 2\left| \frac{\nabla ^{N} (H-|\nu |^2_{\Pten })}{H-|\nu |^2_{\Pten }} + \Pten ^{TN}(\nu , \cdot ) \right| ^2 + |\mathring{\mathrm{II}}+ \mathring{\Pten} ^{TN}|^2\\
	&\ \ + \int_{N_t}\left( 16 \pi (\mu -|J|_g)+ \left( H- \Pten(\nu ,\nu ) \right) \mathrm{tr}_{N_t} \Pten \right)   \\
	&\geq 0.
	\end{split}
\end{align} 

\begin{thm}\label{thm-weak-mono}
	Let $(M^3, g, \kten)$ be a complete, connected, asymptotically flat initial data set without boundary. Let $E_0$ be a precompact open set so that $\partial E_0$ is a connected $C ^{2, \alpha }$-smooth outermost past apparent horizon. Then there exists a weak solution $(N_t) _{0 \leq t< \infty}$ of the $\Pten$-IMCF with initial condition $E_0$ so that $m_{\Pten}(N_t)$ is monotone increasing, whenever the DEC (\ref{eq:DEC}) and the $2$-convexity condition $\mathbf{P} \geq 0$ hold.
\end{thm}

\begin{proof}
	The proof is the same as in \cite[Section 6]{Dong26SigmaIMCF} with necessary modifications.
Using the same notations as in \cite[Section 6]{Dong26SigmaIMCF}, we can replace all $\sigma _i$ by $\mathbf{P}$, and accordingly use the fact that when $t$ is not a jump time, $\tilde{N}_t ^{i} \to \tilde{N}_t$ in $C^{1,\alpha }$ and $|\tilde{\nu }_i|^2_{\mathbf{P}} \to |\tilde{\nu }|^2_{\mathbf{P}}$ in $C^{\alpha }$ on $\tilde{N}_t$. Then the only places where $\sigma ^{i} = |h| g$ has been used is the convergence of the term $\phi (\tilde{\nabla }_{\tilde{\nu }_i} \sigma ^{i})(\tilde{\nu }_i, \tilde{\nu }_i)$ and thus the arguments towards the final monotonicity formula. We indicate necessary modifications as follows.


	Recall that we have the same identity of \cite[Equation (30)]{Dong26SigmaIMCF}:
\begin{align}\label{ddt-identity-i}
	\begin{split}
		&\frac{d}{d t} \int_{\tilde{N}^{i}_t} \phi (\tilde{H}_{i}- |\tilde{\nu}_{i} |^2_{\mathbf{P} })^2 \\							&= \int_{\tilde{N}^{i}_t} (\tilde{H}_{i}-|\tilde{\nu}_i |^2_{ \mathbf{P}}) \tilde{\nabla} _{\tilde{\nu}_i }\phi -2 \langle \tilde{\nabla} ^{N} \phi , \frac{\tilde{\nabla} ^{N}(\tilde{H}_i-|\tilde{\nu}_i |^2_{\mathbf{P}})}{\tilde{H}_i- |\tilde{\nu}_i |^2_{\mathbf{P} }} \rangle \\
		&\ \ - 2 \int_{\tilde{N}^{i}_t}\phi \left(   \frac{|\tilde{\nabla} ^{N} (\tilde{H}_i-|\tilde{\nu}_i |^2_{ \mathbf{P}})|^2}{(\tilde{H}_i-|\tilde{\nu}_i |^2_{\mathbf{P} })^2} + \left( \tilde{\mathrm{Ric}}(\tilde{\nu}_i ,\tilde{\nu}_i ) +|\tilde{\mathrm{II}}|^2 \right) \right) \\
		&\ \ -2 \int_{\tilde{N}^{i}_t} \phi (\tilde{\nabla} _{\tilde{\nu}_i } \mathbf{P}) (\tilde{\nu}_i , \tilde{\nu}_i ) - 4 \int_{\tilde{N}^i_t} \phi \frac{\mathbf{P}( \tilde{\nu}_i , \tilde{\nabla} ^{N}(\tilde{H}_i- |\tilde{\nu}_i |^2_{\mathbf{P} }) )}{\tilde{H}_i- |\tilde{\nu}_i |^2_{\mathbf{P}}} + \int_{\tilde{N}^{i}_t}\phi \tilde{H}_i(\tilde{H}_i-|\tilde{\nu}_i |^2_{\mathbf{P} }).
	\end{split}
\end{align}

Similar to the identity (\ref{nabla-P-id}), we rewrite the term
\begin{align*}
	 (\tilde{\nabla }_{\tilde{\nu }_i} \mathbf{P})(\tilde{\nu }_i, \tilde{\nu }_i) = (\mathrm{div}_{\tilde{g}} \mathbf{P})(\tilde{\nu }_i) - \mathrm{div}_{\tilde{N}^{i}_t}(\mathbf{P}^{TN}(\tilde{\nu }_i, \cdot ) ) - H_{\tilde{N}^{i}_{t}} \mathbf{P}(\tilde{\nu }_i, \tilde{\nu }_i) + \langle \tilde{\mathrm{II}}, \mathbf{P}^{TN} \rangle.
\end{align*} 
Set $\tilde{\nu }_i ^{T} := \tilde{\nu }_i - \langle \tilde{\nu }_i, e_{z} \rangle e_{z}$. We already know that $\sup_{\tilde{N}^{i}_t \cap (M \times \mathrm{supp} \phi )}|\langle \tilde{\nu }_i, e_{z} \rangle| \to 0$ for a.e. $0\leq t \leq T$. So $\tilde{\nu }_i ^{T} \to \tilde{\nu } $ a.e. for a.e. $0\leq t \leq T$. Since $\tilde{g}$ and $\mathbf{P}$ are vertical invariant, we have 
\begin{align*}
	\phi (\mathrm{div}_{\tilde{g}} \mathbf{P})(\tilde{\nu }_i) = \phi (\mathrm{div}_{g} \mathbf{P})(\tilde{\nu }_i ^{T}).
\end{align*} 
By the bounded convergence theorem, for each $0 \leq r <s$, we have
\begin{align*}
	\int_{r}^{s} \int_{\tilde{N}^{i}_{t}} \phi (\mathrm{div}_{\tilde{g}} \mathbf{P})(\tilde{\nu }_i) \to \int_{r}^{s} \int_{\tilde{N}_{t}}\phi (\mathrm{div}_{g} \mathbf{P})(\tilde{\nu }).
\end{align*} 
Using integration by parts, similarly, since $\mathbf{P}(e_{z}, \cdot )=0$, we have
\begin{align*}
	2 \int_{r}^{s}\int_{\tilde{N}^{i}_{t}} \phi \mathrm{div}_{\tilde{N}^{i}_{t}}(\mathbf{P}^{TN}(\tilde{\nu }_i, \cdot ) ) &= 2 \int_{r}^{s}\int_{\tilde{N}^{i}_{t}} \phi ' \langle \tilde{\nu }_i, e_{z} \rangle \mathbf{P}(\tilde{\nu }_{i}, \tilde{\nu }_i ) \to 0.
\end{align*} 
The remaining terms can be controlled exactly as in \cite[Section 6]{Dong26SigmaIMCF}, together with the algebraic identities used to prove (\ref{B'-ineq}). 
\end{proof}

\begin{thebibliography}{99}
\bibitem[Dong26]{Dong26SigmaIMCF} Conghan Dong. \newblock The {$\sigma$}-inverse mean curvature flow and the generalized {Penrose} conjecture, 2026. \newblock arXiv:2605.26504v1.
\end{thebibliography}
\end{document}
